企業や研究機関、知的財産部門では、競合分析、技術ランドスケープ把握、研究開発テーマ探索、ライセンスや共同研究候補の検討など、戦略的意思決定のために特許情報が用いられている。特許文書には、発明が解決しようとする課題、具体的な解決手段、対象物、効果が記述されているため、既存分野の技術動向を把握するだけでなく、他分野で用いられている技術手法や解決手段を探索するための情報源としても利用できる。特に、ある分野で発展した技術手法が別分野の類似課題に対する解決手段となる可能性があるため、異分野間の技術候補探索は、研究開発の方向性や、専門家が確認すべき技術候補を検討するうえで重要である。例えば、医療・生物画像解析で提案され広く用いられてきた U-Net は、後に製造欠陥検出などの非医療画像分野にも応用が広がっている。このような異分野への技術応用候補を、後年の出現を待って確認するだけでなく、特許文書から早期に抽出できれば、研究開発テーマ探索や専門家による技術候補検討を支援できる可能性がある。

一方で、異分野間の技術候補を探索する場合、専門家は自分の専門分野の課題や解決手段には詳しいものの、別分野で類似する課題がどのような語彙、対象物、用途で記述され、どのような解決手段が用いられているかまでは把握しにくい。この問題に対して、特許探索で広く用いられるキーワード検索や分類コードによる絞り込みは、事前に想定した語句や分類に依存しやすく、異分野で異なる表現を用いる候補を見落とす可能性がある。一方で、文書単位の類似度に基づいて候補を集めても、候補が課題、解決手段、対象物、用途のいずれの近さに基づくのかを説明しにくい。したがって、異分野の技術手法・解決手段を専門家確認候補として提示するには、候補がどの課題に対応し、どの解決手段を持つのかを示すだけでなく、その候補が対応する課題の類似性に基づいて抽出されたのか、解決手段の類似性または差分に基づいて抽出されたのか、あるいは対象物や用途の一致に基づいて抽出されたのかを説明できる必要がある。

そこで本研究では、検索式およびCPC・CPCサブクラスに基づいて参照分野と対象分野の特許集合を定義し、異分野間で比較する対象を明確にする。次に、各特許から課題文脈、解決手段文脈、技術手段、対象物、用途分野を抽出し、課題、解決手段、分野依存情報を分けて扱う。課題文脈および解決手段文脈はPatentSBERTaによりベクトル化し、単純な語句一致ではなく、文脈に基づく意味的近さとして比較する。そのうえで、参照分野の特許と対象分野の特許との間で課題文脈が意味的に近く、参照分野では用いられている一方で対象分野では過去時点で低出現または未使用である技術手法・解決手段を専門家確認候補として抽出する。さらに、課題文脈と解決手段文脈の双方が類似する候補を近接解決策型、課題文脈は類似するが解決手段文脈が異なる候補を異質解決策型として整理し、候補提示の根拠を説明できるようにする。

提案する枠組みを確認するため、本研究では二つの実験を行った。実験 1 では、医療画像解析分野で先行して用いられた U-Net を対象に、製造欠陥検出分野への後年出現を参照した後ろ向き評価を行った。その結果、U-Net は過去期間の全体ランキングで上位 10 件に含まれ、具体的手法名に限定したランキングでは 2 位となった。実験 2 では、洗浄・除去系技術から半導体洗浄への異質解決策型候補抽出を補助的に分析した。46 件のユニークファミリーを LLM 支援付き著者確認評価により確認した結果、ノイズを含む一方で、一部の候補は半導体洗浄課題と接続しうる専門家確認候補として解釈可能であった。これらの結果は、課題文脈と解決手段文脈を分けて扱うことで、全文類似とは異なる観点から専門家確認候補を提示できる可能性を示している。

以上より、本研究は、特許文書を課題文脈、解決手段文脈、分野依存情報に分けて扱うことで、全文類似のみでは説明しにくいクロスドメイン技術候補を、専門家確認候補として整理・提示できる可能性を示した。本研究は、特許活用の中でも、技術ランドスケープ分析や研究開発テーマ探索において、異分野技術候補を絞り込むための予備的スクリーニング手法として位置づけられる。ただし、本研究は技術導入可能性や将来の技術移転を証明するものではない。
